
\[
\begin{array}{ll}
\mathrm{i)}& \mathbb R F_k(\underline{L}^{\bullet})\text{ \uline{est} }\mathcal C_0
\text{\uline{-parfait}};\\[2mm]
\mathrm{ii)}& \text{\uline{Le morphisme canonique} }(2.5.3):
\Projz\circ\mathbb R F_k(\underline{L}^{\bullet})\longrightarrow \underline{L}^{\bullet}
\text{ \uline{est un isomorphisme de} }D^+(\mathcal A).
\end{array}
\]

La notion de complexe $\mathcal C_0$-parfait envisagée ici est locale, c'est-à-dire qu'il s'agit de perfection relativement à la catégorie fibrée sur $S$ définie par $\mathcal C_0$. Cette notion de perfection a bien un sens, car $\mathcal C_0$ vérifie les axiomes de (I 4.0) : seule la stabilité par noyau d'épimorphisme n'est pas immédiate ; elle résulte de 1.5.1.

Comme $S$ est quasi-compact, il suffit de vérifier les deux assertions de 2.6 localement sur $S$, de sorte qu'on peut supposer $S$ affine. Le faisceau inversible $\underline{O}_X(1)$ est alors ample, et on dispose de résolutions globales sur $X$ (II 2.2.9). Il existe donc un complexe $\underline{L}'{}^{\bullet}$ de $D_{\mathrm{parf}}(\mathcal A)$, strictement parfait, et un isomorphisme dans $D_{\mathrm{parf}}(\mathcal A)$ : $\underline{L}^{\bullet}\xleftarrow{\sim}\underline{L}'{}^{\bullet}$, ce qui permet de supposer $\underline{L}^{\bullet}$ strictement parfait. Or nous avons vu en 2.4 que, pour $k$ assez grand, les $\underline{L}^i$, qui sont tous nuls sauf un nombre fini, sont $F_k$-acycliques; on a par conséquent~:
\[
  \mathbb R F_k(\underline{L}^{\bullet})=F_k(\underline{L}^{\bullet})
\]
si $k$ est assez grand. Or d'après 1.3 et 1.6, on peut trouver $k_0$ tel que pour $k\geq k_0$ la condition précédente soit remplie, et que les $F_k(\underline{L}^i)$ soient des $\underline{B}$-Modules $\mathcal C_0$-parfaits pour tout $i$. Le complexe $F_k(\underline{L}^{\bullet})$ est alors $\mathcal C_0$-parfait, d'où le i).

Si $k$ est choisi tel que précédemment, l'homomorphisme (2.5.3) est alors l'homomorphisme~:
